\rulesection{6.0 Command}

Command points (CPs) are received during the Command Phase and determine units' movement allowance.

\rulesubsection{6.1 Command Points}

Determine movement allowance as follows:
\begin{itemize}
\item Basic movement allowances for each unit type are listed on the Movement Costs Chart.
\item During the Command Phase, supreme commanders and cavalry commanders receive CPs and broadcast CPs to subordinates.
\item Each unit may use some proportion of its basic movement allowance. The number of CPs received by the senior leader in each hex determines the percentage for all units in the hex.
\end{itemize}

\rulesubsection{6.2 Stacking Effect}

At the instant a unit moves, its movement allowance is modified for units that moved out of the hex previously in that turn.

\rulesubsection{6.3 Receiving Command Points}

During his Command Phase a player determines how many CPs the supreme commander and cavalry commander have for the current turn. Roll once for the supreme commander and once for the cavalry commander. Modify the die roll for line of communications and consult the Command Point Table. The result is the number of CPs received that turn.

\rulesubsection{6.4 Broadcasting}

In broadcast, a commander attempts to give a subordinate the same CPs he has received. The supreme commander and cavalry commander may broadcast the CPs they receive to other friendly leaders not in the same hex.

Commanders never broadcast to units in the same hex; their own command ratings determine movement allowance there. Each commander may broadcast only to subordinates. The supreme commander may not broadcast to the cavalry commander.

Only the senior leader in each hex receives CPs by broadcast. The number received determines movement allowances for every unit in the hex, even if the hex contains both infantry and cavalry.

A commander always attempts to broadcast his full number of CPs and may attempt to broadcast to every hex containing subordinates.

\textbf{6.41 Procedure.} Count the distance in cavalry movement points from the broadcasting commander to the subordinate. Locate the corresponding column on the Command Broadcast Table, roll one die, apply any modifiers for line of communications, and cross-reference the result.

\textbf{6.42 Results.} Results of \texttt{-1}, \texttt{-2}, or \texttt{-3} mean that this many CPs did not get through. If the number lost is greater than or equal to the number broadcast, the subordinate receives no CPs. A result of \texttt{Full} means the subordinate receives the full number of CPs the broadcasting commander received.

\textbf{6.43 Broadcast Range.} Broadcast range is counted in cavalry movement points, not hexes. Any distance traced along a friendly intact rail line counts as one hex regardless of distance or terrain. Movement costs for zones of control must be paid unless a friendly combat unit occupies the hex. There is no maximum broadcast distance. The broadcast path may not be traced through enemy units.

\rulesubsection{6.5 Command Points and Movement}

\textbf{6.51} Compare the supreme commander's CPs to his command rating and express the ratio as a percentage. This is the percentage of his basic movement allowance he may expend.

\textbf{6.52} All units stacked with the supreme commander use the same percentage of their movement allowance, even if they do not remain with him during movement.

\textbf{6.53} For all other units, use the senior leader in each hex. Modify his command rating if necessary and compare broadcast CPs to that rating. Every leader in the hex receives the same percentage of movement allowance as the senior leader.

\textbf{6.54} If a wagon is stacked with a leader and the leader has a movement allowance of 6 or less, the wagon's movement allowance is 6. If it is not stacked with a leader, use the wagon's own command rating.

\textbf{6.55 Minimum Movement Allowance.} A unit that remains more than six hexes from all enemy units during its entire move may always expend 50\% of its maximum movement allowance, regardless of CPs. This does not entitle it to entrench or destroy or repair.

\rulesubsection{6.6 Modifier to Command Rating}

If three or more combat units are together in a hex and no commander is present, add 1 to the senior leader's command rating when determining movement allowances or commitment.

\rulesubsection{6.7 Restrictions}

Leaders receiving more CPs than their command rating may not save them, nor does this increase movement allowance. Leaders may only increase movement allowance through force march. A unit may not move if its hex receives 0 CPs, except as provided by destination markers or the minimum movement allowance rule.

\rulesubsection{6.8 Line of Communications}

When using the Command Point Table, a player must check whether the commander being rolled for has a line of communications (LOC). When using the Command Broadcast Table, he must check whether the broadcasting commander has an LOC.

\textbf{6.81} A supreme or cavalry commander traces an LOC to a friendly depot. The LOC may be traced over at most four non-road hexes to a road or rail hex, and then any distance along a continuous path of road or rail hexes to a friendly depot.

\textbf{6.82} A commander traces an LOC to his supreme commander the same way the supreme commander traces one to a depot. The supreme commander must also be able to trace an LOC himself.

\textbf{6.83} If a supreme or cavalry commander cannot trace an LOC, there is a \texttt{-1} DRM on both the Command Table and the Command Broadcast Table.

\textbf{6.84} An LOC may not be traced through an enemy unit. Zones of control do not affect LOC. No portion of an LOC may be traced through an unbridged river, and the road or rail portion may not cross a destroyed bridge or destroyed rail hex.
